\section{Methods\label{chapter3}}

This chapter describes the experimental framework used to compare Echo-State Networks (via the pyReCo library) and LSTM networks for chaotic time-series prediction. Section~\ref{sec:framework} provides an overview of the experimental design. Section~\ref{sec:data_prep} describes dataset generation and preprocessing. Section~\ref{sec:model_config} details the model configurations and parameter budget matching strategy. Section~\ref{sec:metrics} defines the evaluation metrics. Section~\ref{sec:stat_methods} describes the statistical analysis methods.

% ============================================================================
\subsection{Experimental Framework Overview}
\label{sec:framework}
% ============================================================================

The experimental framework is designed to provide a fair comparison between PyReCo (Echo-State Networks) and LSTM networks under controlled conditions. The key design principles are:

\begin{enumerate}
    \item \textbf{Matched parameter budgets}: Both architectures are evaluated at three parameter budget levels (small, medium, large) with approximately equal total parameter counts.
    \item \textbf{Multiple datasets}: Three benchmark systems spanning synthetic chaotic (Lorenz, Mackey-Glass) and real-world experimental (Santa Fe laser) data.
    \item \textbf{Data efficiency analysis}: Training data sizes are varied systematically from 1\,000 to 10\,000 samples.
    \item \textbf{Statistical validity}: Each experiment is repeated with five random seeds (42, 43, 44, 45, 46), and significance is assessed through hypothesis testing.
    \item \textbf{Comprehensive metrics}: Beyond prediction accuracy, training time, and computational efficiency are measured.
\end{enumerate}

The complete experimental campaign comprises 417 experiments: 147 main experiments (3 datasets $\times$ 3 budgets $\times$ 3 training fractions $\times$ 2 models, with 5 seeds each) and 270 data efficiency experiments (3 datasets $\times$ 6 data lengths $\times$ 3 budgets $\times$ 5 seeds, for both models). All experiments are implemented in Python using PyTorch for LSTM and the pyReCo library for ESN models. Reproducibility is ensured through fixed random seeds and environment locking.

% ============================================================================
\subsection{Dataset Preparation}
\label{sec:data_prep}
% ============================================================================

% --------------------------------------------------------------------------
\subsubsection{Data Generation}
\label{sec:data_gen}
% --------------------------------------------------------------------------

\textbf{Lorenz system.} Trajectories are generated by integrating the Lorenz equations (Eqs.~\ref{eq:lorenz_x}--\ref{eq:lorenz_z}) using \texttt{scipy.integrate.solve\_ivp} with a time step of $\Delta t = 0.01$. The canonical parameters are used: $\sigma = 10$, $\rho = 28$, $\beta = 8/3$, with initial condition $\mathbf{x}_0 = (1, 1, 1)$. An initial transient is discarded to allow the trajectory to settle onto the attractor. The three-dimensional output $(x, y, z)$ is available, with the $x$-component used as the primary prediction target. The default data length is 15\,000 time steps; for data efficiency experiments, lengths of 1\,000, 2\,000, 3\,000, 5\,000, 7\,000, and 10\,000 are used.

\textbf{Mackey-Glass system.} The delay differential equation (Eq.~\ref{eq:mackeyglass}) is integrated using Euler integration with a step size of $h = 1.0$, initial condition $x_0 = 1.2$, and parameters $\beta_{\text{MG}} = 0.2$, $\gamma = 0.1$, $n = 10$, $\tau = 17$ (chaotic regime). The default data length is 20\,000 time steps. A warmup period is discarded, and the resulting scalar time series is used directly.

\textbf{Santa Fe laser dataset.} The dataset is loaded from the \texttt{reservoirpy} library (\texttt{reservoirpy.datasets.santafe\_laser}). It consists of approximately 10\,000 intensity measurements from an NH\textsubscript{3} far-infrared laser \cite{weigend1994time}.

% --------------------------------------------------------------------------
\subsubsection{Preprocessing}
\label{sec:preprocessing}
% --------------------------------------------------------------------------

All time series are normalized using \texttt{StandardScaler}, which transforms each feature to zero mean and unit variance:
\begin{equation}
    x_{\text{scaled}} = \frac{x - \bar{x}}{s},
    \label{eq:standard_scaler}
\end{equation}
where $\bar{x}$ and $s$ are the mean and standard deviation computed from the \emph{training set only}. The same transformation is applied to validation and test sets to prevent data leakage.

\textbf{Sliding window creation.} The normalized time series is transformed into input-output pairs using a sliding window approach. Given a lookback window of $n_{\text{in}}$ steps and a forecast horizon of $n_{\text{out}}$ steps, each sample consists of an input sequence $[x_{t-n_{\text{in}}+1}, \ldots, x_t]$ and a target $[x_{t+1}, \ldots, x_{t+n_{\text{out}}}]$. The default configuration uses $n_{\text{in}} = 100$ and $n_{\text{out}} = 1$ for single-step prediction. The choice of $n_{\text{in}} = 100$ follows the pyReCo library default and is consistent with values commonly used in the reservoir computing literature for chaotic time-series prediction \cite{jaeger2001echo, lukosevicius2012practical, chattopadhyay2020data, pathak2018model}.

% --------------------------------------------------------------------------
\subsubsection{Train/Validation/Test Split}
\label{sec:split}
% --------------------------------------------------------------------------

The data is split temporally (without shuffling) to respect the sequential nature of time series:
\begin{itemize}
    \item \textbf{Training set}: The first fraction of the data, controlled by the \texttt{train\_frac} parameter (typically 0.7 for main experiments; varied across 0.2, 0.3, 0.5, 0.6, 0.7, 0.8 for data efficiency studies).
    \item \textbf{Validation set}: A fixed fraction of 15\% of the data, used for LSTM early stopping and PyReCo hyperparameter selection via cross-validation.
    \item \textbf{Test set}: The remaining portion (e.g., 15\% when \texttt{train\_frac} = 0.7), used exclusively for final evaluation. No model selection or hyperparameter tuning is performed on the test set.
\end{itemize}

For data efficiency experiments, the training fraction is fixed at 0.7 while the total data length is varied from 1\,000 to 10\,000 samples, allowing the study of how absolute data availability affects model performance.

% ============================================================================
\subsection{Model Configurations}
\label{sec:model_config}
% ============================================================================

% --------------------------------------------------------------------------
\subsubsection{Parameter Budget Matching Strategy}
\label{sec:budget_matching}
% --------------------------------------------------------------------------

A central contribution of this thesis is the comparison of PyReCo and LSTM under matched \emph{total} parameter budgets. Unlike trainable parameters alone, total parameters capture the full computational footprint of the model, including the fixed reservoir weights in ESNs. Three budget levels are defined:

\begin{table}[H]
\centering
\caption{Parameter budget levels and corresponding model configurations.}
\label{tab:budgets}
\begin{tabular}{@{}llrrrr@{}}
\toprule
Budget & Model & Hidden/Nodes & Total Params & Trainable Params & Trainable \% \\
\midrule
\multirow{2}{*}{Small (${\sim}$1K)}
    & PyReCo & 100 nodes & ${\sim}$1\,000 & 100 & ${\sim}$10\% \\
    & LSTM   & 10 hidden & 952 & 952 & 100\% \\
\midrule
\multirow{2}{*}{Medium (${\sim}$10K)}
    & PyReCo & 300 nodes & ${\sim}$10\,000 & 300 & ${\sim}$3\% \\
    & LSTM   & 49 hidden & 10\,052 & 10\,052 & 100\% \\
\midrule
\multirow{2}{*}{Large (${\sim}$50K)}
    & PyReCo & 700 nodes & ${\sim}$50\,000 & 700 & ${\sim}$1.4\% \\
    & LSTM   & 109 hidden & 49\,077 & 49\,077 & 100\% \\
\bottomrule
\end{tabular}
\end{table}

For PyReCo, the number of reservoir nodes $N$ is selected to match the target parameter budget based on the total parameter formula (Eq.~\ref{eq:esn_total}), accounting for the configured density and input fraction. For LSTM, the hidden size $h$ is selected such that $4h(h + d + 1) + h \cdot d_{\text{out}} + d_{\text{out}}$ matches the budget. All readout nodes in the ESN contribute to the output (\texttt{fraction\_output} = 1.0).

This matching strategy means that at the small budget, a 100-node ESN with only 100 trainable readout weights is compared against an LSTM with 952 fully trainable parameters---both having approximately 1\,000 total parameters. The implications of this choice are discussed in Chapter~\ref{chapter5}.

% --------------------------------------------------------------------------
\subsubsection{PyReCo Configuration}
\label{sec:pyreco_config}
% --------------------------------------------------------------------------

The PyReCo model uses the following hyperparameter search grid, yielding 36 combinations per experiment:

\begin{table}[H]
\centering
\caption{PyReCo hyperparameter search grid (36 combinations).}
\label{tab:pyreco_grid}
\begin{tabular}{@{}lll@{}}
\toprule
Hyperparameter & Values & Motivation \\
\midrule
Spectral radius ($\rho$)    & 0.8, 0.9, 1.0 & Memory / stability trade-off \\
Leakage rate ($\alpha$)      & 0.4, 0.5, 0.6 & Dynamics speed \\
Density ($\delta$)           & 0.05, 0.1     & Sparsity of reservoir \\
Input fraction ($f_{\text{in}}$) & 0.5, 0.75  & Input connectivity \\
\bottomrule
\end{tabular}
\end{table}

The best hyperparameter combination is selected using 5-fold time-series cross-validation with forward-chaining (expanding window), which respects the temporal ordering of the data. The model achieving the lowest validation MSE across folds is selected for final evaluation on the test set.

ESN hyperparameter sensitivity has been extensively documented in the literature \cite{lukosevicius2012practical, matzner2022hyperparameter, mwamsojo2024stochastic}. The chosen ranges are centered on commonly recommended values, with the spectral radius range spanning the critical transition near $\rho = 1$.

% --------------------------------------------------------------------------
\subsubsection{LSTM Configuration}
\label{sec:lstm_config}
% --------------------------------------------------------------------------

The LSTM model uses a single-layer architecture with hidden size matched to the parameter budget (Table~\ref{tab:budgets}). The hyperparameter search grid comprises 9 combinations:

\begin{table}[H]
\centering
\caption{LSTM hyperparameter search grid (9 combinations).}
\label{tab:lstm_grid}
\begin{tabular}{@{}lll@{}}
\toprule
Hyperparameter & Values & Motivation \\
\midrule
Learning rate & 0.0005, 0.001, 0.002 & Most critical hyperparameter \cite{greff2017lstm} \\
Dropout       & 0.1, 0.2, 0.3       & Regularization \cite{zaremba2014recurrent} \\
\bottomrule
\end{tabular}
\end{table}

The fixed training parameters are:
\begin{itemize}
    \item Optimizer: Adam \cite{kingma2015adam}
    \item Loss function: Mean Squared Error (MSE)
    \item Maximum epochs: 100
    \item Early stopping: Patience of 10 epochs monitoring validation loss
    \item Batch size: 32
\end{itemize}

The choice of learning rate as the primary tuning parameter follows the finding of Greff et al. \cite{greff2017lstm} that learning rate is ``by far the most important hyperparameter'' for LSTM, with performance dropping significantly if set incorrectly. Dropout rates between 0.1 and 0.3 are recommended for sequence learning tasks \cite{gal2016dropout}.

Both models receive systematic hyperparameter tuning to ensure a fair architectural comparison, though the ESN grid (36 combinations) is larger than the LSTM grid (9 combinations) to accommodate the known higher sensitivity of ESNs to hyperparameter choices \cite{matzner2022hyperparameter}.

% ============================================================================
\subsection{Evaluation Metrics}
\label{sec:metrics}
% ============================================================================

% --------------------------------------------------------------------------
\subsubsection{Prediction Accuracy}
\label{sec:accuracy_metrics}
% --------------------------------------------------------------------------

The following metrics quantify prediction accuracy. Let $\{y_i\}_{i=1}^{n}$ denote the true values and $\{\hat{y}_i\}_{i=1}^{n}$ the predicted values.

\textbf{Mean Squared Error (MSE):}
\begin{equation}
    \text{MSE} = \frac{1}{n} \sum_{i=1}^{n} (y_i - \hat{y}_i)^2.
    \label{eq:mse}
\end{equation}

\textbf{Mean Absolute Error (MAE):}
\begin{equation}
    \text{MAE} = \frac{1}{n} \sum_{i=1}^{n} |y_i - \hat{y}_i|.
    \label{eq:mae}
\end{equation}

\textbf{Root Mean Squared Error (RMSE):}
\begin{equation}
    \text{RMSE} = \sqrt{\text{MSE}}.
    \label{eq:rmse}
\end{equation}

\textbf{Normalized Root Mean Squared Error (NRMSE):}
\begin{equation}
    \text{NRMSE} = \frac{\text{RMSE}}{y_{\max} - y_{\min}}.
    \label{eq:nrmse}
\end{equation}

\textbf{Coefficient of Determination (R\textsuperscript{2}):}
\begin{equation}
    R^2 = 1 - \frac{\sum_{i=1}^{n}(y_i - \hat{y}_i)^2}{\sum_{i=1}^{n}(y_i - \bar{y})^2},
    \label{eq:r2}
\end{equation}
where $\bar{y} = \frac{1}{n}\sum_{i=1}^{n} y_i$. An $R^2$ of 1 indicates perfect prediction; 0 indicates performance no better than predicting the mean; negative values indicate worse-than-mean predictions. $R^2$ is used as the primary comparison metric throughout this thesis because it is scale-invariant and directly interpretable.

% --------------------------------------------------------------------------
\subsubsection{Multi-Step Prediction}
\label{sec:multi_step}
% --------------------------------------------------------------------------

In addition to single-step ($n_{\text{out}} = 1$) prediction, models are evaluated on multi-step forecasting with horizons of 1, 5, 10, 20, and 50 steps. Multi-step predictions are generated in \emph{free-run} (autoregressive) mode: the model's own predictions are fed back as inputs for subsequent steps, without access to ground-truth values. This mode is the most challenging because errors compound with each step.

Multi-step evaluation is critical for chaotic systems because the predictability horizon is fundamentally limited by the positive Lyapunov exponent. The rate at which prediction quality degrades with horizon length provides insight into how well each model captures the underlying dynamics.

% --------------------------------------------------------------------------
\subsubsection{Computational Efficiency}
\label{sec:efficiency_metrics}
% --------------------------------------------------------------------------

Following the methodology outlined in our evaluation protocol, computational efficiency is measured through three separate timing components, reported independently to provide a complete picture:

\begin{enumerate}
    \item \textbf{Training time with optimal parameters (primary metric)}: Wall-clock time for a single training run using the best hyperparameter configuration identified during tuning. For PyReCo, this is the time to construct the reservoir and compute the ridge regression readout. For LSTM, this is the time for gradient-based training with early stopping. This metric represents the fair comparison, as both models are evaluated in their ``best state,'' and is the metric most commonly reported in academic literature.

    \item \textbf{Hyperparameter tuning time (supplementary)}: Wall-clock time for the complete hyperparameter search. For PyReCo, this includes 5-fold cross-validation over 36 hyperparameter combinations ($36 \times 5 = 180$ training runs). For LSTM, this includes training across 9 hyperparameter combinations with a single validation split ($9 \times 1 = 9$ training runs). The tuning process is classified as a preparatory methodology step---performed once per configuration to identify optimal hyperparameters---and is reported separately from the primary training time comparison.

    \item \textbf{Inference time}: Wall-clock time per test sample for generating predictions, measured after training is complete.
\end{enumerate}

\textbf{Training efficiency} is defined as $R^2$ per second of training time (with optimal parameters), capturing the rate at which predictive performance is acquired.

All timing measurements are performed on the same hardware to ensure fair comparison. The rationale for separate reporting is that hyperparameter tuning is performed once per (dataset, budget) configuration, while training with optimal parameters may be repeated many times (e.g., across random seeds or for retraining on new data). In practical deployment, the tuning cost is amortized, making the training-only time the more relevant metric for model selection.

% ============================================================================
\subsection{Statistical Analysis}
\label{sec:stat_methods}
% ============================================================================

All comparisons are evaluated for statistical significance using the following methods:

\textbf{Paired t-test.} For each configuration (dataset, budget, data length), PyReCo and LSTM are trained with the same five random seeds. The paired t-test assesses whether the mean difference in performance (e.g., MSE or $R^2$) between the two models is significantly different from zero:
\begin{equation}
    t = \frac{\bar{d}}{s_d / \sqrt{n}},
    \label{eq:ttest}
\end{equation}
where $\bar{d}$ is the mean of the paired differences, $s_d$ is the standard deviation of the differences, and $n = 5$ is the number of pairs. The null hypothesis (no difference between models) is rejected at significance level $\alpha = 0.05$.

\textbf{Wilcoxon signed-rank test.} As a non-parametric alternative that does not assume normality of the differences \cite{wilcoxon1945individual}, the Wilcoxon signed-rank test is also applied to confirm the results of the parametric t-test.

\textbf{Cohen's d effect size.} To quantify the magnitude of differences beyond statistical significance, Cohen's d is computed \cite{cohen1988statistical}:
\begin{equation}
    d = \frac{\bar{d}}{s_d}.
    \label{eq:cohens_d}
\end{equation}
Effect sizes are interpreted as: $|d| < 0.2$ (negligible), $0.2 \leq |d| < 0.5$ (small), $0.5 \leq |d| < 0.8$ (medium), $|d| \geq 0.8$ (large).

\textbf{95\% Confidence intervals} are reported for all mean differences to quantify the uncertainty of the estimates.

With 54 pairwise comparisons across all configurations, the probability of at least one false positive at $\alpha = 0.05$ is substantial. Results are therefore interpreted considering both the individual p-values and the overall pattern of significance across configurations, rather than relying on any single comparison.
